\section{Odd primes and the finite-place lower bound}
\label{sec:odd-primes}

Throughout this section assume that $G\in\mathbb Q$, so that the columns
$F_j$ of Section~\ref{sec:foundations} are rational.  We regard each $F_j$
as a column in $\mathbb Q^n$.  For an odd prime $p$, reduction modulo $p$
always means reduction of a member of $\mathbb Z_p$; in particular, every
congruence below includes the assertion that its two sides are locally
integral.

The raw matrix has $n+q$ columns, whereas the filtered determinant has
size $n$.  At the large primes treated first, each raw column has at most
two powers of $p$ in its denominator.  We shall make invertible column
changes over $\mathbb Z_p$ and bound how many columns can still have
denominator $p^2$ or $p$.  Such bounds control every full minor of the
raw matrix, and Cauchy--Binet then transfers them to the integer column
filter.  The first reduction identifies residues after multiplication
by $p^2$; for $p>H/2$ a second reduction identifies the remaining
residues after multiplication by $p$.

\subsection{Digit reduction of the rational moments}

We first prove the reductions needed at primes
\[
  2\sqrt H<p\le H.
\]
We may suppose that $p$ does not divide the denominator of $G$: as $N$
tends to infinity, every fixed denominator prime eventually lies below
$2\sqrt H$.  Put
\[
 E(t)=(1-t^2)^{(p-1)/2}=\sum_{d=0}^{p-1}E_dt^d
 \quad\text{in }\mathbb F_p[t],\qquad
 \epsilon=(-1)^{(p-1)/2}.
\]
We set $E_d=0$ outside $0\le d<p$.  Within this range,
\[
 E_d=\begin{cases}c_{d/2}\pmod p,&d\text{ even},\\0,&d\text{ odd},\end{cases}
 \qquad E_{p-1-d}=\epsilon E_d.
\]
The first identity follows from
$(-1)^k\binom{(p-1)/2}{k}\equiv4^{-k}\binom{2k}{k}$; the second follows
by reversing the coefficients of $E$.

\begin{lemma}[Digit reductions]\label{lem:odd-digit-reductions}
For $0\le i,j<H$, let
\[
 \ell=j\bmod p,\quad d=(j-i-1)\bmod p,
 \qquad
 i'=\frac{i+1+d-\ell}{p}-1,\quad j'=\left\lfloor\frac jp\right\rfloor,
 \qquad 0\le\ell,d<p.
\]
Then $i'\ge-1$, and, with the convention
$M^0(-1,j')=k^+_{j'+1}$,
\begin{equation}\label{eq:digit-layers}
 \begin{split}
 p^2M^0(i,j)&\equiv E_dM^0(i',j')\pmod p,\\
 p^2Z^0(i,j)&\equiv
 \begin{cases}
 Z^0(\lfloor i/p\rfloor,j'),&i\equiv j\pmod p,\\
 0,&i\not\equiv j\pmod p.
 \end{cases}
 \end{split}
\end{equation}
\end{lemma}

\begin{proof}
We first give the parity and carry calculation for the factors $H_z^*$
in the explicit formulas of Section~\ref{sec:foundations}.  Write
$z=Pp+r$, $0\le r<p$, $0\le z<H$.  Since $H<p^2/4$, all factors at
the reduced indices $P$ are $p$-adic units.  If $z$ is even, then
\begin{equation}\label{eq:odd-even-H}
 \begin{cases}
 H_z^*\equiv H_P^*c_{r/2}\pmod p,&r\text{ even},\\
 v_p(H_z^*)=1,&r\text{ odd}.
 \end{cases}
\end{equation}
Indeed, when $r$ is even, $P$ is even and the base-$p$ digits of $z/2$
are $P/2,r/2$.  Expanding $(1+x)^z$ modulo $p$ by
$(1+x)^p=1+x^p$ gives the first assertion, including the factor $4^{-z/2}$.
When $r$ is odd, $P$ is odd and the low digit of $z/2$ is $(p+r)/2$.
The factorial formula for $\binom z{z/2}$ has exactly one carry:
\[
 \left\lfloor\frac zp\right\rfloor
       -2\left\lfloor\frac{z/2}{p}\right\rfloor=1,
\]
and there is no contribution from $p^2$.  Thus, for positive even $z$,
$v_p(zH_z^*)$ is either zero or one.  It is one precisely when $r=0$
or $r$ is odd.

For odd $z$ one has
\begin{equation}\label{eq:odd-odd-H}
 \frac1{H_z^*}\in\mathbb Z_p,\qquad
 pH_z^*\equiv
 \begin{cases}
 \epsilon c_{r/2}H_P^*,&r\text{ even},\\
 0,&r\text{ odd}.
 \end{cases}
\end{equation}
If $r$ is odd, both $z$ and $c_{(z-1)/2}$ are units by the same
digit calculation.  If $r$ is even, $P$ is odd.  At $z=Pp$, digit
expansion gives
\[
 c_{(Pp-1)/2}\equiv c_{(P-1)/2}c_{(p-1)/2}
       =\epsilon c_{(P-1)/2},
 \qquad pH_{Pp}^*\equiv\epsilon H_P^*.
\]
The recurrence $H_{z+2}^*/H_z^*=(z+1)/(z+2)$, applied through
$r=0,2,\ldots,p-1$, now proves the nonzero case of
Equation~\eqref{eq:odd-odd-H}.  Its denominators on this interval are
units.  Since $m_{z-1}=2H_z^*$ for odd $z$ and is zero for even $z$,
these formulas imply
\begin{equation}\label{eq:odd-moment-digit}
 pm_{z-1}\equiv E_{p-1-r}m_{P-1}\pmod p.
\end{equation}
This includes $z=0$, using $m_{-1}=0$.

We next reduce the two starting arrays.  For $u=Pp+r$, $0\le u<H$,
we claim
\begin{equation}\label{eq:odd-starting-digits}
 p^2k^-_u\equiv E_{p-1-r}k^-_P,
 \qquad
 p^2k^+_{u+1}\equiv E_r k^+_{P+1}\pmod p.
\end{equation}
For odd $u$, write the first formula as
\[
 p^2 k^-_u=(pH_u^*)
           \sum_{\substack{1\le z\le u\\z\text{ odd}}}\frac{2p}{z}.
\]
Only $z=mp$, with $m$ odd, survives in the sum.  If $r$ is odd the
first factor vanishes modulo $p$.  If $r$ is even, $P$ is odd and
Equation~\eqref{eq:odd-odd-H} gives
$\epsilon E_rH_P^*\sum_{m\le P,\ m\text{ odd}}2/m
 =E_{p-1-r}k^-_P$, as required.

For even $u$ the quantities $p/(zH_z^*)$, for positive even $z\le u$,
are integral.  Thus
\[
 p^2k^-_u=H_u^*
     \sum_{\substack{1\le z\le u\\z\text{ even}}}
                  2\left(\frac{p}{zH_z^*}\right)^2
\]
is integral, and is zero modulo $p$ when $r$ is odd.  If $r$ is even,
then $P$ is even.  The nonzero summands form precisely the complete
blocks
\[
 z=mp-\lambda,\qquad
 m=2,4,\ldots,P,\qquad \lambda=0,2,\ldots,p-1.
\]
To see completeness, the indices with $z\bmod p$ odd lie immediately
below an even multiple $mp$, while that multiple supplies $\lambda=0$.
The upper limit $u=Pp+r$ contains the whole block ending at $Pp$,
and the next such block begins at $(P+1)p+1>u$.  No partial block is
present.  The recurrence, started at $mp$ and followed downwards, gives
\[
 \frac{p}{(mp-\lambda)H_{mp-\lambda}^*}
       \equiv\frac{c_{\lambda/2}}{mH_m^*}\pmod p.
\]
For example, each downward step from $z$ to $z-2$ multiplies this
quantity by $(z-1)/(z-2)$; these multipliers give successively
$(2j-1)/(2j)$ modulo $p$.  Finally, with $s=(p-1)/2$,
\[
 \sum_{j=0}^{s}c_j^2
  \equiv\sum_{j=0}^{s}\binom{s}{j}^2
  =\binom{2s}{s}\equiv(-1)^s=\epsilon\pmod p.
\]
The equality is the coefficient identity obtained from
$(1+x)^s(1+x)^s$.  Multiplying the block sums by
$H_u^*\equiv E_rH_P^*$ proves the first congruence of
Equation~\eqref{eq:odd-starting-digits} also for even $u$.

For the second congruence, suppose first that $u$ is even.  In
\[
 p^2k^+_{u+1}
   =H_u^*\sum_{\substack{1\le z\le u\\z\equiv u\ (2)}}
                     \frac{2p^2}{z^2H_z^*},
\]
every term with $p\nmid z$ vanishes modulo $p$, because
$v_p(H_z^*)\le1$.  At $z=mp$, necessarily $m$ is even and
$H_{mp}^*\equiv H_m^*$.  The result is $E_r k^+_{P+1}$ when $r$ is
even, and zero when $r$ is odd.  For odd $u$ instead write the expression
as
\[
 (pH_u^*)\sum_{\substack{1\le z\le u\\z\text{ odd}}}
                         \frac{2p}{z^2H_z^*}.
\]
The reciprocal $1/H_z^*$ is integral.  Again only multiples $z=mp$
can survive; there
$pH_{mp}^*\equiv\epsilon H_m^*$, and the two factors of $\epsilon$
cancel.  If $r$ is odd the prefactor is zero.  This proves the second
congruence in every case and also proves the local integrality asserted
in Equation~\eqref{eq:odd-starting-digits}.

If $i\ge j$, put $u=i-j=Pp+r$.  Then $d=p-1-r$ and $i'=j'+P$.
The diagonal reduction is
\[
 M^0(i,j)=k^-_u-\sum_{k=1}^{j}\frac{m_{u+k-1}}k.
\]
After multiplication by $p^2$, terms with $p\nmid k$ vanish by
Equation~\eqref{eq:odd-moment-digit}.  At $k=mp$ the surviving term is
$E_{p-1-r}m_{P+m-1}/m$.  The first starting congruence therefore gives
\[
 p^2M^0(i,j)\equiv E_{p-1-r}
       \left(k^-_P-\sum_{m=1}^{j'}\frac{m_{P+m-1}}m\right)
       =E_dM^0(i',j').
\]
If $i<j$, put $u=j-i-1=Pp+r$, so $d=r$ and $i'=j'-P-1$.  Now
\[
 M^0(i,j)=k^+_{u+1}-\sum_{k=u+2}^{j}\frac{m_{k-u-2}}k.
\]
For a multiple $k=mp$ in the sum,
\[
 k-u-1=(m-P-1)p+(p-1-r).
\]
Equation~\eqref{eq:odd-moment-digit} gives the reduced summand
$E_r m_{m-P-2}/m$.  A possible first multiple with $m=P+1$ contributes
zero, since its reduced moment is $m_{-1}$.  Thus the reduced sum is
over $m=P+2,\ldots,j'$, as in the formula for $M^0(i',j')$.
The expression for $i'$ is also
$i'=\lfloor(i-\ell)/p\rfloor$, so $i'\ge-1$.  When $i'=-1$,
$j'=P$, the sum is empty and the starting term is exactly
$k^+_{j'+1}=M^0(-1,j')$.  This proves the first assertion of
Equation~\eqref{eq:digit-layers}, including its boundary convention.
The reduced array indices are less than $H/p<p/4$; even the boundary
index $j'+1$ is less than $p$.  Thus the explicit rational formulas
show that all quantities on its right side are integral at $p$.

For $Z^0$, write $I=\lfloor i/p\rfloor$ and $J=\lfloor j/p\rfloor$.
Because $i,j<p^2$, the harmonic sums satisfy
\[
 pB_i^{(1)}\equiv B_I^{(1)},\qquad
 p^2B_i^{(2)}\equiv B_I^{(2)}\pmod p.
\]
The diagonal formula follows at once.  Off the diagonal, if
$i\not\equiv j\pmod p$, the denominator $i-j$ is a unit and
$p^2Z^0(i,j)$ is zero modulo $p$.  If the residues agree,
$i-j=p(I-J)$ with $I-J$ a unit; substituting the first harmonic
congruence gives $Z^0(I,J)$.  This also proves integrality of
$p^2Z^0(i,j)$.
\end{proof}

\subsection{The leading layer and paired raw columns}

We now gather the moment reductions into column vectors.  This will
identify pairs of raw columns whose sum has at most one power of $p$ in
its denominator.  For a fixed residue $0\le\ell<p$, define column vectors over
$\mathbb F_p$ by
\[
 P'_u=[t^{(u+1)p+\ell}]\,tP(t)E(t)\quad(u\ge-1),
 \qquad
 D'_u=[t^{up+\ell}]\,D(t)\quad(u\ge0).
\]
Here $P,D$ denote the vectors of all row polynomials; the dependence
of these extracted vectors on $\ell$ is understood.  Their supports
are finite.  Lemma~\ref{lem:odd-digit-reductions}, gathered by coefficient
residue, gives
\begin{equation}\label{eq:column-layer}
 X_{kp+\ell}:=p^2F_{kp+\ell}\bmod p
   =\sum_{u\ge-1}P'_uM^0(u,k)
          -\frac32\sum_{u\ge0}D'_uZ^0(u,k).
\end{equation}
For a term $[t^i]P$, its unique residue $d$ satisfies
$i+1+d=(i'+1)p+\ell$, so that it contributes exactly to $P'_{i'}$,
including $i'=-1$.  The $Z^0$ reduction survives only when
$i=up+\ell$, which is precisely the extraction defining $D'_u$.
The coefficient of $G$ in $F_j$ is a sum of integral multiples
of $c_{(i-j)/2}$.  These coefficients are integral at every odd prime,
so this term disappears after multiplication by $p^2$.  Every raw
column therefore lies in $p^{-2}\mathbb Z_p^n$.

If $\ell\ge H-2p$, the degree bounds
$\deg(tPE)\le H+p-1$ and $\deg D<H$ show that $P'_u,D'_u$ vanish
for $u>1$.  Put
\[
 V_\ell=2P'_1-\frac32D'_1,
 \qquad U_\ell=2P'_{-1}-\frac32D'_0.
\]
The needed small values of the reduced arrays are
\[
 \begin{array}{c|rrr}
 u&-1&0&1\\\hline
 M^0(u,0)&0&0&2\\
 M^0(u,1)&2&0&-2
 \end{array}
 \qquad
 \begin{array}{c|rr}
 u&0&1\\\hline
 Z^0(u,0)&0&1\\
 Z^0(u,1)&1&-1
 \end{array}.
\]
Consequently, whenever the indicated columns exist,
\begin{equation}\label{eq:odd-paired-layer}
 X_\ell=V_\ell,\qquad X_{p+\ell}=U_\ell-V_\ell.
\end{equation}
Since $tP$ and $D$ have no terms below degree $A+1$,
$U_\ell=0$ for $\ell\le A$.

We record explicitly how changes in the full column pool will be used.
Let $T$ be the $n$ by $(n+q)$ matrix of all raw columns.  If
$Q\in\operatorname{GL}_{n+q}(\mathbb Z_p)$ and every full minor of
$TQ$ has valuation at least $-D$, the same holds for every full minor
of $T=(TQ)Q^{-1}$ by Cauchy--Binet.  The coefficients in this expansion
are minors of the integral matrix $Q^{-1}$.  Applying Cauchy--Binet
once more to the integer matrix defining the prescribed filter gives
$v_p(\Delta_N)\ge-D$.  This transfer applies even when the columns or
their leading residues are linearly dependent.

For $2\sqrt H<p\le H/2$, use each pair of columns with
\[
 \max(b,H-2p)\le\ell<\min(p,L-p,A).
\]
Their number is exactly
\[
 d_0=\bigl(\min(p,L-p,A)-\max(b,H-2p)\bigr)_+,
 \qquad y_+=\max(y,0).
\]
Replacing its high column by the sum of the pair is an integral
elementary column operation with integral inverse.  The pair sum lies
in $p^{-1}\mathbb Z_p^n$ by Equation~\eqref{eq:odd-paired-layer}; all
other columns still lie in $p^{-2}\mathbb Z_p^n$.  These pairs are
disjoint, since $\ell<p$.  Any selection of $n$ columns from the pool
of $n+q$ contains at least $(d_0-q)_+$ of the improved columns.
The transfer just explained gives
\begin{equation}\label{eq:odd-lower-pair-range}
 v_p(\Delta_N)\ge-2n+(d_0-q)_+.
\end{equation}

\subsection{The central layer and integral column elimination}

For $p>H/2$, the next lemma identifies the residues of central columns
after multiplication by $p$.  It expresses them in terms of the same
vectors $V_\ell$ that occur in the leading residues of noncentral
columns.  The elimination proof will retain noncentral columns and use
$p$ times those columns to cancel the corresponding $V_\ell$ terms.
\begin{lemma}[Central layer]\label{lem:odd-central-layer}
Suppose $H/2<p\le H$, and write
$K=L-p$ and $J=H-p$.  For every central column, by which we mean
\[
 b\le j<\min(p,L),\qquad j\ge J,
\]
one has
\begin{equation}\label{eq:central-layer}
 pF_j\in\mathbb Z_p^n,\qquad
 pF_j\equiv-\sum_{0\le\ell<J}\frac{V_\ell}{j-\ell}\pmod p.
\end{equation}
All denominators $j-\ell$ occurring here are units at $p$.
\end{lemma}

\begin{proof}
For every contributing row degree $0\le i<H$,
$i-j\le H-1-J=p-1$ and $j-i<p$.  Hence $|i-j|<p$, and the
starting value $k^-_{i-j}$ or $k^+_{j-i}$ in the diagonal reduction
is integral at $p$.  In either case that reduction can be written as
the starting value minus
\[
 \sum_{0\le\ell<\min(i,j)}\frac{m_{i-\ell-1}}{j-\ell}.
\]
Here $1\le j-\ell<p$.  By Equation~\eqref{eq:odd-moment-digit}, a
nonzero reduction after multiplication by $p$ requires
$p\le i-\ell<2p$, and then
\[
 pm_{i-\ell-1}\equiv2E_{2p+\ell-1-i}.
\]
Such an index necessarily has $0\le\ell<J$.  Conversely its
coefficient can be collected over all $i$ with no truncation, since
$j\ge J>\ell$ and $i\ge p+\ell>\ell$.  Thus
\[
 pM^0(P,j)\equiv-\sum_{0\le\ell<J}\frac{2P'_1}{j-\ell}.
\]
The extracted vector $P'_1$ in each summand is the one belonging to
that residue $\ell$.
For $Z^0(i,j)$, the harmonic sums have a pole only if
$i=p+\ell$ with $0\le\ell<J$.  The denominator $i-j$ is then a
unit: equality modulo $p$ would require $j=\ell<J$.  Therefore
\[
 pZ^0(p+\ell,j)\equiv\frac1{\ell-j}.
\]
All other $Z^0(i,j)$ are integral.  Combining these two calculations
with the factor $-3/2$ proves Equation~\eqref{eq:central-layer}; the
$G$ term is integral here as before.
\end{proof}

We now carry out these cancellations in the full column pool.  In the
resulting pool, $R$ will count the retained columns with possible
denominator $p^2$, and $S$ will bound the number of other columns with
possible denominator $p$; every remaining column will be integral.
The proof first uses the $V_\ell$ supplied by retained noncentral
columns to modify the central columns.  We then use the rank of their
remaining scaled residues to bound the number of central columns that
can still have denominator $p$ after an invertible column change.

\begin{lemma}[Elimination over the local integers]
\label{lem:odd-local-elimination}
For $H/2<p\le H$, define
\begin{align*}
 R&=K_++(K-A)_++(J-\max(b,K))_+,\\
 S&=(\min(A,K)-b)_+
        +(\min(b,J)-\max(0,K))_+.
\end{align*}
Then every full raw minor, and hence the filtered determinant, satisfies
\begin{equation}\label{eq:odd-two-layer-bound}
 v_p(\Delta_N)\ge-\min(2n,n+R,2R+S).
\end{equation}
\end{lemma}

\begin{proof}
The noncentral columns are precisely the high columns
$F_{p+\ell}$ for $0\le\ell<K$ and the low columns $F_\ell$ for
$b\le\ell<J$.  Empty ranges are allowed.  Indeed $p>b$, $J<p$,
and $K<J$.  For
\[
 b\le\ell<\min(A,K)
\]
retain the low column and replace the high column by
$F_{p+\ell}+F_\ell$.  Let
\[
 B=(\min(A,K)-b)_+
\]
be the number of these pair sums.  Each lies in $p^{-1}\mathbb Z_p^n$.
Retain every other noncentral column as a column in
$p^{-2}\mathbb Z_p^n$, whether or not its leading residue depends on
the others.  The number of columns so retained is
\begin{equation}\label{eq:odd-retained-count}
 R_0=K_++(J-b)_+-B.
\end{equation}
It equals the stated $R$.  To verify the identity, split $K$ into
$K\le0$, $0<K\le b$, $b<K\le A$, and $K>A$; use $K<J$ in the
last two cases.  The resulting formulas are respectively
$(J-b)_+$, $K+(J-b)_+$, $J$, and $K+J-A$.

These retained columns furnish a representative of every vector
$V_\ell$ with $0\le\ell<J$ except possibly those in
\[
 \mathcal E=\{\ell:\max(0,K)\le\ell<\min(b,J)\}.
\]
For $b\le\ell<J$, the retained low column has leading residue
$p^2F_\ell\equiv V_\ell$.  For $0\le\ell<\min(b,K)$, the
unpaired high column has leading residue $-V_\ell$, because
$\ell<b<A$ makes $U_\ell=0$.  These two ranges exhaust the complement
of $\mathcal E$.  Denote the size of $\mathcal E$ by
\[
 e=|\mathcal E|=(\min(b,J)-\max(0,K))_+.
\]

For each represented residue choose a retained column $Y_\ell$ and
a sign $\sigma_\ell\in\{1,-1\}$ such that
$p^2Y_\ell\equiv\sigma_\ell V_\ell$.  For every central $j$, choose
$a_{\ell j}\in\mathbb Z_p$ reducing to
$\sigma_\ell/(j-\ell)$ and replace that column by
\[
 C_j=F_j+p\sum_{\substack{0\le\ell<J\\\ell\notin\mathcal E}}
                                      a_{\ell j}Y_\ell.
\]
This is a shear of the whole column pool with integral coefficients
and inverse obtained by changing the signs of the added coefficients.
It preserves every retained column and every pair sum.  The resulting
central columns lie in $p^{-1}\mathbb Z_p^n$ and satisfy
\[
 pC_j\equiv-\sum_{\ell\in\mathcal E}\frac{V_\ell}{j-\ell}
       \pmod p.
\]
This step does not require any information about the next coefficient
of $Y_\ell$.  Explicitly, if
$Y_\ell=p^{-2}y_0+p^{-1}y_1+y_2$ with integral lifts $y_0,y_1,y_2$,
then $pY_\ell=p^{-1}y_0+y_1+py_2$; its unknown second coefficient is
already integral.

Let $c$ be the number of central columns.  Their scaled residues
define a linear map
\[
 \varphi:\mathbb F_p^c\longrightarrow\mathbb F_p^n,
 \qquad (a_j)_j\longmapsto\sum_j a_j(pC_j\bmod p).
\]
Its image is contained in the span of the $e$ exceptional vectors,
so $d:=\operatorname{rank}\varphi\le\min(c,e)$.  Choose a basis of
$\mathbb F_p^c$ whose last $c-d$ vectors form a basis of
$\ker\varphi$.  Lift its basis matrix arbitrarily to a matrix
$W\in\operatorname{Mat}_{c}(\mathbb Z_p)$.  Its determinant is a
unit, so $W\in\operatorname{GL}_c(\mathbb Z_p)$ and the adjugate
formula gives $W^{-1}\in\operatorname{Mat}_{c}(\mathbb Z_p)$.
Applying this change to the central columns leaves at most $d$
columns in $p^{-1}\mathbb Z_p^n$; the other $c-d$ are integral,
because their scaled residues vanish and they already belonged to
$p^{-1}\mathbb Z_p^n$.  If $c=0$ this step is the empty identity.

The resulting full pool consequently contains $R$ columns with
possible denominator $p^2$, at most $B+d\le B+e=S$ other columns
with possible denominator $p$, and integral remaining columns.
No noncentral column was discarded, nor was its denominator reduced
on the strength of a leading linear dependence.  In a full minor,
if $r$ columns come from the first group and $s$ from the second,
the loss is at most $2r+s$.  The inequalities
\[
 2r+s\le2n,\qquad 2r+s\le n+R,\qquad 2r+s\le2R+S
\]
hold since $r+s\le n$, $r\le R$, and $s\le S$.  All changes used
above and their inverses are integral, so the Cauchy--Binet transfer
proves Equation~\eqref{eq:odd-two-layer-bound} for all original full
minors and for $\Delta_N$.
\end{proof}

The omission of the optional pair at $\ell=A$ is harmless: its
physical column is included in $R$.  The preceding proof also covers
$K\le0$ and $J\le b$.  In particular, if $L\le p\le H$, there
are no high or low noncentral columns, $R=0$, and $S=H-p$; all raw
columns are central.  At the formal endpoint $p=H$, all of them are
integral.

\subsection{Summing the prime losses}

For a nonzero rational determinant, the valuations weighted by
$\log p$ sum to its ordinary logarithmic absolute value.  We now
combine the odd-prime estimates with the bound at $2$ to obtain the
lower bound used in the final comparison.

\begin{proposition}[Finite-place bound]\label{prop:finite-lower}
Assume $G\in\mathbb Q$.  Along any sequence of positive integers
$N\to\infty$ for which $\Delta_N\ne0$, one has
\begin{equation}\label{eq:finite-lower-final}
 \liminf_{N\to\infty}
 \left(\frac{\log|\Delta_N|}{n^2}-\frac12\log2\right)
 \ge -\frac{8609}{4608}
       -\left(\frac12+\frac{505}{4608}\right)\log2
 >-2.29084.
\end{equation}
\end{proposition}

\begin{proof}
We first compute the complete loss function.  Put $x=p/N$.  For
$x\le65/2$, Equation~\eqref{eq:odd-lower-pair-range} has
\[
 \frac{d_0}{N}
  =\bigl(\min(x,59-x,19)-\max(7,65-2x)\bigr)_+.
\]
For $0\le x\le23$ this is zero; for $23\le x\le29$ it is
$2x-46$; and for $29\le x\le65/2$ it is $12$.  Only the part
exceeding $q/N=4$ improves a full minor.  For $x>65/2$, the formulas
of Lemma~\ref{lem:odd-local-elimination} give the following values;
the table retains the intermediate breakpoint $52$ at which the
counting formula changes:
\[
\begin{array}{c|cc}
 x\text{ range}&R/N&S/N\\\hline
 {[65/2,40]}&105-2x&12\\
 {[40,52]}&65-x&52-x\\
 {[52,58]}&117-2x&x-52\\
 {[58,59]}&59-x&6\\
 {[59,65]}&0&65-x
\end{array}.
\]
The entries agree at their common endpoints.  Taking
$\min(96,48+R/N,2R/N+S/N)$ yields the additional breakpoint $69/2$.
Thus for every relevant prime $p>2\sqrt H$ the valuation is bounded
below by $-Nd(p/N)$, where $d$ is continuous, is zero for $x\ge65$,
and has the following exact pieces:
\begin{equation}\label{eq:odd-loss-table}
\begin{array}{c|c|c|c}
 x\text{ range}&d(x)&
       \text{endpoint values}&\displaystyle\int_{\text{range}}d(x)\,dx\\\hline
 {[0,25]}&96&96,96&2400\\
 {[25,29]}&146-2x&96,88&368\\
 {[29,65/2]}&88&88,88&308\\
 {[65/2,69/2]}&153-2x&88,84&172\\
 {[69/2,40]}&222-4x&84,62&803/2\\
 {[40,58]}&182-3x&62,8&630\\
 {[58,59]}&124-2x&8,6&7\\
 {[59,65]}&65-x&6,0&18
\end{array}.
\end{equation}
In particular,
\begin{equation}\label{eq:odd-loss-area}
 \int_0^{65}d(x)\,dx=\frac{8609}{2}.
\end{equation}

For completeness, the omitted small odd primes cost only $o(n^2)$
in the logarithm of the determinant.  At any odd prime, the factorial
formula gives
\[
 0\le v_p\binom{2l}{l}\le1+\frac{\log H}{\log p}
 \qquad(2l\le H).
\]
Each factorial-floor difference is zero or one.  The displayed
formulas for $H_z^*,m_i,k_d^\pm,Z^0$, and diagonal reduction then
give, with an absolute constant $C_0$ and with
$\operatorname{den}G$ denoting the positive reduced denominator of $G$,
\[
 v_p(F_j(r))\ge
   -C_0\left(1+\frac{\log H}{\log p}\right)-v_p(\operatorname{den}G).
\]
For instance, each summand in $k^-_u$ uses at most two powers of
$z$ and two powers of $H_z^*$ in its denominator; no additional
loss arises from summing.  All row and filter coefficients are
integers.  Expanding a determinant therefore multiplies this entry
bound by at most $n$, and summing it over $p\le2\sqrt H$ with
weights $\log p$ gives
\[
 O\bigl(n\sqrt H\log H+n\log(\operatorname{den}G)\bigr)
       =o(n^2).
\]
This estimate uses only that the number of such primes is at most
$2\sqrt H$.  For sufficiently large $N$, every prime $p>H$ is
integral throughout the calculation: factorials and degree
denominators introduce only prime factors at most $H$, and the fixed
denominator of $G$ has no prime factor above $H$.  Thus these primes
give no negative contribution.

We use the classical prime number theorem in the form
\[
 \theta(y):=\sum_{p\le y}\log p\sim y
       \qquad(y\to\infty);
\]
see \cite[Step~VI, p.~707]{Zagier1997}, which presents Newman's analytic
proof. The weighted consequence needed here follows directly:
\begin{equation}\label{eq:odd-weighted-primes}
 \frac1N\sum_{p\le65N}d(p/N)\log p
       \longrightarrow\int_0^{65}d(x)\,dx.
\end{equation}
Here is a direct justification of the weight.  For a fixed partition
$0=x_0<\cdots<x_m=65$, the prime number theorem gives
\[
 \frac{\theta(Nx_j)-\theta(Nx_{j-1})}{N}
       \longrightarrow x_j-x_{j-1}.
\]
Upper and lower step functions formed from the supremum and infimum
of $d$ on each interval bound the weighted sum.  Their limiting
difference tends to zero as the mesh decreases, since $d$ is
uniformly continuous.  Individual partition endpoints change the
sum by at most $O(\log N/N)$.  This proves
Equation~\eqref{eq:odd-weighted-primes}.  The prime $2$ and the
primes below $2\sqrt H$ may be removed from that sum at cost $o(1)$,
using $\|d\|_\infty=96$ and the elementary bound
$\theta(2\sqrt H)\le2\sqrt H\log(2\sqrt H)$.

Combining these estimates with Equation~\eqref{eq:odd-loss-area}
and $n^2=2304N^2$ gives
\[
 \sum_{p\text{ odd}}v_p(\Delta_N)\log p
       \ge-\frac{8609}{4608}n^2-o(n^2).
\]
Proposition~\ref{prop:two-adic}, with $\delta=10/48$, gives
\[
 v_2(\Delta_N)\ge
  -\left(\frac\delta2+\frac{\delta^2}{8}\right)n^2-o(n^2)
  =-\frac{505}{4608}n^2-o(n^2).
\]
For a nonzero rational number its numerator and denominator
factorizations give the exact identity
\[
 \log|\Delta_N|=\sum_pv_p(\Delta_N)\log p.
\]
Subtracting $\tfrac12\log2$ after dividing by $n^2$ proves the
non-strict inequality in Equation~\eqref{eq:finite-lower-final}.

The final numerical comparison also has a short rational check.
The identity
\[
 \log2=2\sum_{j=0}^{\infty}\frac{3^{-(2j+1)}}{2j+1}
\]
and its positive geometric tail give
\[
 \log2\le
 2\sum_{j=0}^{5}\frac{3^{-(2j+1)}}{2j+1}
       +\frac{2\,3^{-13}}{13(1-3^{-2})}
 <\frac{693149}{10^6}.
\]
Consequently the lower constant is strictly larger than
\[
 -\frac{8609}{4608}-\frac{2809}{4608}\frac{693149}{10^6}
 =-\frac{10556055541}{4608000000}>-2.29084,
\]
which completes the proof.
\end{proof}
